\section{Reglamentos de industria y estándares internacionales en sistemas de ordeño automatizado y datos pecuarios}

\subsection{Justificación}

En la minería de datos aplicada a la ganadería de precisión, las variables capturadas por los sensores de un robot de ordeño (flujo de leche, conductividad eléctrica por cuarto, recuento de células somáticas, frecuencias de visita e identificadores de radiofrecuencia RFID) no son datos abstractos: están estrictamente reguladas por estándares internacionales y códigos de calidad comercial de la industria láctea.

Si un equipo de ciencia de datos ignora estos estándares:

\begin{enumerate}
  \item Construirá modelos de Machine Learning sobre métricas biológicamente inválidas o mal calculadas (e.g. estimar producciones diarias sin ponderar los intervalos simétricos del ordeño voluntario).
  \item Vulnerará acuerdos de confidencialidad comercial y normas de inocuidad alimentaria que rigen la relación entre el establo y las plantas pasteurizadoras receptoras de la leche.
\end{enumerate}

\subsection{Reglamentos de la industria y cámaras de comercio}

\begin{table}[H]
  \centering
  \small
  \begin{tabularx}{\textwidth}{>{\raggedright\arraybackslash}p{0.24\textwidth}>{\raggedright\arraybackslash}X>{\raggedright\arraybackslash}X}
    \toprule
    Reglamentos de la Industria Láctea Nacional & & Impacto en reto CAETEC \\
    \midrule
    CANILEC (Cámara nacional de industriales de la leche) &
    \textbullet~Estándares de recepción de leche cruda\newline \textbullet~Confidencialidad de calidad higiénica\newline \textbullet~Cláusulas de penalización comercial &
    Protección de Métricas de Tanque y Salud (Conductividad MDI y RCS como confidencialidad de negocio) \\
    \addlinespace
    FEMELECHE (Federación Mexicana de Lechería) &
    \textbullet~Protección Patrimonial del Productor\newline \textbullet~Soberanía y Titularidad de Datos de Hato\newline \textbullet~Costos y Transparencia en Cadena Láctea &
    Garantía de no divulgación de productividad (curvas de lactancia y acuerdos de compraventa) \\
    \addlinespace
    SENASICA / SADER (Servicio nacional de sanidad, inocuidad y calidad agroalimentaria) &
    \textbullet~Manual de Buenas Prácticas Pecuarias (BPPL-01)\newline \textbullet~Trazabilidad y Bitácoras Sanitarias Inmutables &
    Obligación de trazabilidad auditable (acceso a registros de salud animal) \\
    \bottomrule
  \end{tabularx}
\end{table}

\subsubsection{Cámara nacional de industriales de la leche}

\textbf{Recurso oficial:} Estándares de calidad, inocuidad y evaluación de la conformidad en la recepción de leche cruda (CANILEC, compendio técnico en concordancia con NOM-155-SCFI-2012 y NOM-243-SSA1-2010).

\textbf{Fundamentos y regularización del sector:} Agrupa las principales plantas procesadoras del país. Regula las tablas de bonificación y penalización aplicadas a los establos lecheros con base en variables que aparecen en los archivos del reto

\begin{enumerate}
  \item Calidad Fisiológica y Composición: Porcentaje de grasa, proteína y sólidos no grasos (estimados por el software del VMS).
  \item Calidad Higiénica y Sanitaria: Recuento de Células Somáticas (RCS x 1000 en \texttt{patadas\_180725.csv}). Un establo con RCS $>$ 400,000 cel/mL sufre castigos económicos severos o el rechazo total del volumen del silo
\end{enumerate}

\textbf{Justificación legal para proyecto:} Los reportes de calidad higiénica de un establo son estrictamente confidenciales, la exposición o filtración de los mismos afectará la reputación comercial. Por lo tanto, las variables de salud de ubre deben tratarse bajo esquemas de secreto industrial y disociarse del nombre y razón social de la explotación.

\subsubsection{Federación Mexicana de Lechería (FEMELECHE)}

\textbf{Recurso oficial:} Lineamientos para la Competitividad, Trazabilidad y Salvaguarda de la Información Pecuaria en la Lechería Familiar y Especializada (FEMELECHE, 2022).

\textbf{Fundamento y regulación en el sector:} FEMELECHE representa los intereses de los productores frente a la cadena de valor y el gobierno federal. Establece que la información técnica y de costos de los establos pertenece al productor ganadero y no puede ser utilizada por proveedores tecnológicos (DeLaval) ni intermediarios para distorsionar precios de mercado

\textbf{Justificación para el proyecto:} Como estudiantes e investigadores nos encontramos obligados a respetar la soberanía sobre la información ganadera estipulada por FEMELECHE: los registros de rendimiento diario (como por ordeño o el volumen global acumulado) deben ser evaluados exclusivamente con fines científicos institucionales, bloqueando la cesión de estas métricas hacia terceros u organizaciones lucrativas sin la autorización explícita y previa del socio formador

\subsubsection{SENASICA / SADER: Buenas Prácticas Pecuarias en Leche Bovina}

\textbf{Recurso oficial:} Manual de buenas prácticas pecuarias en unidades de producción de leche bovina (Clave Técnica: BPPL-01, Servicio Nacional de Sanidad, Inocuidad y Calidad Agroalimentaria - SENASICA).

\textbf{Fundamento y regulación en sector:} Exige que todo establo tecnificado cuente con un sistema fidedigno e inalterable de registro individual de cada vaca (nacimiento, partos, tratamientos estéticos, secados y retiros de leche por antibióticos).

\textbf{Justificación para el proyecto:} La columna \emph{Días desde el último tratamiento de salud} y \emph{Días desde el último tratamiento de vacunación} en \texttt{reporte\_180725.csv} obedece directamente a la exigencia de SENASICA de respetar los tiempos de retiro de leche contaminada.

\subsection{Estándares Internacionales y órganos regulatorios mundiales}

El ICAR es la máxima autoridad técnica global (con sede en Roma) que homologa la recolección, cálculo y certificación de datos fenotípicos y genéticos en rumiantes.

\subsubsection{ICAR Device Certification y Estándares RFID (ISO 11784 / ISO 11785)}

\textbf{Recurso oficial:} ICAR Conformance Testing of RFID Devices for Animal Identification

\textbf{Justificación para el reto:} estructura del campo Transponder ISO para Id de 15 dígitos:

\begin{itemize}
  \item Primeros 3 dígitos: Código de país ISO 3166 (484 para México) o código de fabricante aprobado por ICAR (900 al 999, como el 999XXXXXXXXXXXX del archivo de CAETEC).
  \item Siguientes 12 dígitos: Identificador numérico nacional irrepetible del animal.
  \item Conclusión legal: La norma internacional demuestra que este número es un identificador primario universal. Si el equipo no lo enmascara criptográficamente, cualquier auditor que aplique el estándar ICAR identificará la granja de origen.
\end{itemize}

\subsubsection{IDF/FIL (International Dairy Federation / Federación Internacional de Lechería)}

La IDF es la organización científica internacional de referencia ante la FAO y el Codex Alimentarius para la cadena láctea.

\minisec{Boletines oficiales IDF sobre mastitis y células somáticas}

\textbf{Recurso oficial:}

\begin{itemize}
  \item Bulletin of the IDF No. 448/2011: Suggested Interpretation of Mastitis Terminology.
  \item Bulletin of the IDF No. 466/2013: Guidelines for the use and interpretation of bovine milk somatic cell counts (SCC) in the dairy industry.
  \item Bulletin of the IDF No. 515/2022: Guidelines for defining quarter and udder health status and cured clinical and subclinical mastitis cases.
\end{itemize}

\textbf{Justificación:}

\begin{enumerate}
  \item Umbral de células somáticas: La IDF fija internacionalmente que un valor de $\mathit{RCS} \le 200{,}000\ \mathit{cel/mL}$ indica un cuarto mamario fisiológicamente sano. valores $> 200{,}000\ \mathit{cel/mL}$ clasifican inequívocamente a la vaca con mastitis subclínica
  \item Índice de Conductividad por cuarto (MDI): La IDF establece que las variaciones relativas de conductividad eléctrica entre cuartos mamarios (cuarto anterior izquierdo/derecho y cuarto posterior izquierdo/derecho, variables “DI, DD, TI, TD”) superiores al 15--20\% constituyen una alerta temprana de inflamación epitelial antes de que aparezcan grumos en la leche
\end{enumerate}

\minisec{Guia IDF sobre bienestar animal en ordeño robótico}

\textbf{Recurso oficial:} The IDF Guide on Animal Welfare in Dairy Production (Boletín FIL-IDF No.504/2020).

\textbf{Justificación Objetivo enfoque predicción de colas:} La IDF dictamina que el tiempo máximo que una vaca debe permanecer parada en áreas de espera de hormigón antes del ordeño no debe superar los 60 a 90 minutos diarios acumulados.

\textbf{Impacto en el reto:} Si el robot VMS entra en congestión crítica ($>$ 90\% de utilización desordenada) y las vacas esperan más de 2 horas en la manga:

\begin{itemize}
  \item Se incrementa la incidencia de patologías podales (cojeras y lesiones en pezuñas, lo que explica la correlación directa con la tabla de patadas por dolor durante el ordeño).
  \item Se reduce el tiempo de reposo en cubículos y rumia, deprimiendo la producción lechera en hasta un 15\%.
\end{itemize}

Esto justifica formalmente ante el evaluador por qué predecir la saturación de la cola del robot es un objetivo de bienestar animal y salud respaldado por la IDF

\subsection{Artículos científicos indexados que respaldan uso de los estándares}

\subsubsection*{Paper 1 Integración de Sensores y Estándares de Detección Sanitaria}

\begin{itemize}
  \item Referencia: Rutten, C. J., Velthuis, A. G. J., Steeneveld, W., \& Hogeveen, H. (2013). Invited review: Sensors to support health management on dairy farms. Journal of Dairy Science, 96(4), 1928–1952.
  \item DOI: 10.3168/jds.2012-6107 \url{https://doi.org/10.3168/jds.2012-6107}
  \item Revista / Cuartil: Journal of Dairy Science (Elsevier / American Dairy Science Association), JCR Q1.
  \item Aporte para el proyecto: Revisa exhaustivamente cómo los sensores de los robots de ordeño (conductividad eléctrica, color, flujo y células somáticas) se validan frente a los estándares de la IDF y la directriz ICAR. Demuestra que los algoritmos de Learning deben entrenarse sobre las desviaciones relativas entre cuartos ($\Delta$MDI) y no sobre valores absolutos, fundamentando el preprocesamiento de \texttt{patadas\_180725.csv}.
\end{itemize}

\subsubsection*{Paper 2 Tráfico de Ganado, Bienestar y Teoría de Colas en Sistemas Robóticos}

\begin{itemize}
  \item Referencia: Gargiulo, J. I., Eastwood, C. R., Garcia, S. C., \& Lyons, N. A. (2018). Dairy cow traffic in automated milking systems: Challenges and opportunities. Journal of Dairy Science, 101(12), 10609–10626.
  \item DOI: 10.3168/jds.2018-14634 \url{https://doi.org/10.3168/jds.2018-14634}
  \item Revista / Cuartil: Journal of Dairy Science, JCR Q1.
  \item Aporte para el proyecto: Analiza la formación de colas de espera en los boxes de ordeño voluntario y el impacto del tiempo ocioso (idling time) del robot. Fundamenta con rigor biológico y matemático por qué la predicción de congestión con 2 horas de anticipación protege el bienestar del animal conforme a las directrices de la IDF y FA
\end{itemize}

\subsection*{Referencias}

\begin{enumerate}
  \item ICAR (International Committee for Animal Recording). (2023). ICAR Guidelines for Dairy Cattle Milk Recording: Section 2 (Milk Recording) and Section 11 (Automated Milking Systems). ICAR, Roma, Italia. \url{https://www.icar.org/Guidelines/}
  \item IDF (International Dairy Federation). (2020). The IDF Guide on Animal Welfare in Dairy Production. Bulletin of the International Dairy Federation No. 504/2020, Bruselas, Bélgica. \url{https://fil-idf.org/}
  \item IDF (International Dairy Federation). (2022). Guidelines for defining quarter and udder health status and cured clinical and subclinical mastitis cases. Bulletin of the International Dairy Federation No. 515/2022, Bruselas, Bélgica.
  \item CANILEC (Cámara Nacional de Industriales de la Leche). (2020). Compendio de Normas y Parámetros Técnicos de Recepción de Leche Cruda. México, D.F. \url{https://canilec.org.mx/}
  \item SENASICA – SADER. (2019). Manual de Buenas Prácticas Pecuarias en Unidades de Producción de Leche Bovina (BPPL-01). Dirección General de Inocuidad Agroalimentaria, Acuícola y Pesquera, México.
  \item Rutten, C. J., Velthuis, A. G. J., Steeneveld, W., \& Hogeveen, H. (2013). Invited review: Sensors to support health management on dairy farms. Journal of Dairy Science, 96(4), 1928–1952. DOI: 10.3168/jds.2012-6107.
  \item Gargiulo, J. I., Eastwood, C. R., Garcia, S. C., \& Lyons, N. A. (2018). Dairy cow traffic in automated milking systems: Challenges and opportunities. Journal of Dairy Science 101(12), 10609–10626. DOI: 10.3168/jds.2018-14634.
\end{enumerate}
